\section{Introduction}
\label{sec:introduction}

Information retrieval in the food domain poses challenges beyond standard text matching, because relevance depends on structured culinary knowledge---ingredient composition, category hierarchy, preparation style, and flavor compatibility---rather than lexical overlap alone \cite{manning2008introduction, robertson2009probabilistic, karpukhin2020dense, min2022applications}. This difficulty is amplified in Vietnamese cuisine, where accent variation, bilingual naming, and regional diversity interact with a large ingredient vocabulary \cite{nguyen2020phobert}.

Dense retrieval pipelines improve semantic matching over lexical methods by encoding queries and documents in a shared embedding space \cite{karpukhin2020dense, gao2023rag}, yet dense retrieval alone exhibits three recurrent failure modes in the food domain \cite{barnett2024seven, edge2024local}. First, \emph{class-level queries} such as \textit{``mon protein thuc vat''} (plant-protein dishes) cannot be resolved without a hierarchical mechanism to expand the query into descendant ingredients. Second, \emph{constrained ingredient substitution} (e.g., replacing beef under a vegetarian constraint) requires typed relations and class-based filtering absent from flat embeddings \cite{shirai2021identifying, fatemi2023learning}. Third, \emph{related-dish recommendation} via ingredient overlap cannot distinguish within-class substitutions from cross-class ones, because flat Jaccard treats both equivalently \cite{park2021flavorgraph}.

An ontology addresses these limitations by encoding \emph{subclass-of} relations, named relations (\emph{substitutes}, \emph{flavorComplements}, \emph{cookedBy}), and inference rules that enable class-level query resolution through ancestor--descendant traversal \cite{gruber1993translation, dooley2018foodon}. Prior work confirms that explicit concept-level structure improves retrieval precision in specialized domains \cite{vallet2005ontology, fernandez2011semantic, castells2007adaptation}, and recent studies show that coupling retrieval with knowledge-graph traversal complements neural retrieval on compositional queries \cite{sun2023think, he2024g}.

This paper proposes an ontology-enhanced retrieval framework for Vietnamese food IR. Our contributions are fourfold:
\begin{itemize}
  \item A Vietnamese food ontology over ${\sim}$10{,}000 dishes, comprising a four-level ingredient hierarchy, a two-axis dish taxonomy (\emph{byType} $\times$ \emph{byMethod}), and seven named relations (Section~\ref{sec:framework}).
  \item Integration of the ontology into a dense-retrieval pipeline at three injection points: query expansion, constrained substitution reasoning, and hierarchy-aware similarity scoring.
  \item Three evaluation tasks isolating ontology contribution: class-based retrieval (200 queries), constrained substitution (100 cases with LLM-judge ground truth), and related-dish recommendation (200 anchors with diverse candidates and component ablation), with inter-annotator agreement reported (Section~\ref{sec:tasks}).
  \item A four-system comparison with paired significance tests, providing direct evidence for the independent contribution of hierarchy, typed relations, and inference rules.
\end{itemize}

Section~\ref{sec:related_work} reviews related work. Section~\ref{sec:framework} presents the ontology-enhanced framework. Section~\ref{sec:tasks} defines the evaluation tasks. Section~\ref{sec:experiments} reports results and ablation. Section~\ref{sec:conclusion} concludes.
